% --- Archon physics formalization source begin ---
% archon:physics
% archon:covers PhyXMiniProblems/problem_phyx_mini_0282.lean
% archon:source-report reports/phyx_mini/problem_phyx_mini_0282.source.json

\chapter{Physics problem phyx\_mini\_0282}
\label{ch:PhyXMiniProblems_problem_phyx_mini_0282}

\paragraph{Problem source.}
\#\# Physical scenario

The physical pendulum in figure has two possible pivot points \$A\$ and \$B\$. Point \$A\$ has a fixed position but \$B\$ is adjustable along the length of the pendulum as indicated by the scaling. When suspended from \$A\$, the pendulum has a period of \$T = 1.80\$ s. The pendulum is then suspended from \$B\$, which is moved until the pendulum again has that period.

\#\# Question

What is the distance \$L\$ between \$A\$ and \$B\$?

\#\# Answer choices

- A: A: 0.768 m

- B: B: 0.786 m

- C: C: 0.826 m

- D: D: 0.804 m

\#\# Figure caption (auxiliary; use the image as primary evidence)

The image depicts a wedge-like structure shaded in orange, with several elements annotated on it. Here’s a detailed description:

- The shape is irregular and seems to have a pivot at the top.
- Point "A" is marked at the top left of the shape, slightly protruding from the contour.
- A line from point "A" extends downwards diagonally to point "B".
- Between points "A" and "B", there is a midpoint labeled "com".
- The length of the line from "A" to "B" is marked as "L".
- A dashed vertical line runs through the figure, passing near point "A".
- There are double-headed arrows at the bottom of the dashed line, suggesting rotational movement around the vertical axis.
- "B" is closer to the bottom of the shape with vertical lines extending horizontally across the shape, resembling a series of steps or markers.

The arrangement suggests a pendulum-like setup pivoting around point "A" with the center of mass ("com") indicated.

\#\# Dataset metadata

Resonance and Harmonics; Multi-Formula Reasoning, Physical Model Grounding Reasoning

\paragraph{Recorded answer/context.}
D: D: 0.804 m

\paragraph{Figure/image path.}
/root/proposal\_for\_physic/science-mango/phyx\_mini\_run/phyx\_data/test\_image/282.png

\paragraph{Formalization target.}
create a compiling Lean file with sorry bodies at `PhyXMiniProblems/problem\_phyx\_mini\_0282.lean`. The Lean declarations must preserve the physical quantities, dimensions or dimensional roles, figure labels, governing-law hypotheses, and final relation expressed by this problem.
Use Mathlib/Physlib names found through LeanExplore where available. If a physics API is missing, introduce faithful local abstractions rather than scalar placeholder aliases.

\begin{theorem}[Physics formalization target]
\label{thm:physics:phyx_mini_0282:target}
\lean{PhyXMiniProblems.ProblemPhyXMini0282.pivotSeparation_is_answer_D}
\uses{def:physics:phyx-mini-0282:phyxminiproblems-problemphyxmini0282-lengthinmeters, def:physics:phyx-mini-0282:phyxminiproblems-problemphyxmini0282-reversiblephysicalpendulum, def:physics:phyx-mini-0282:phyxminiproblems-problemphyxmini0282-matchesprimaryfigure, def:physics:phyx-mini-0282:phyxminiproblems-problemphyxmini0282-hasreversiblependulumproblemdata, def:physics:phyx-mini-0282:phyxminiproblems-problemphyxmini0282-satisfiesreversiblephysicalpendulumlaws, def:physics:phyx-mini-0282:phyxminiproblems-problemphyxmini0282-equivalentsimplependulumlengthmeters, def:physics:phyx-mini-0282:phyxminiproblems-problemphyxmini0282-answerlengthmeters, def:physics:phyx-mini-0282:phyxminiproblems-problemphyxmini0282-recordedanswerchoice, def:physics:phyx-mini-0282:phyxminiproblems-problemphyxmini0282-isuniqueclosestanswerchoice}
The assigned autoformalize agent should translate this physics problem into Lean declarations in the covered file, with theorem and lemma proof bodies written as `by sorry`.
\end{theorem}
\begin{proof}
This is an autoformalization task, not a proof task. Produce statements that can later be proved by the physics prover without weakening the physical model.
\end{proof}

% --- Archon declaration topology begin ---
\section{Declaration topology}
\begin{definition}[Length Quantity]
  \label{def:physics:phyx-mini-0282:phyxminiproblems-problemphyxmini0282-lengthquantity}
  \lean{PhyXMiniProblems.ProblemPhyXMini0282.LengthQuantity}
  A nonnegative physical length, independent of the chosen unit system
\end{definition}
\begin{proof}
  This is the declared model definition of Length Quantity.
\end{proof}
\begin{definition}[Time Quantity]
  \label{def:physics:phyx-mini-0282:phyxminiproblems-problemphyxmini0282-timequantity}
  \lean{PhyXMiniProblems.ProblemPhyXMini0282.TimeQuantity}
  A nonnegative physical duration
\end{definition}
\begin{proof}
  This is the declared model definition of Time Quantity.
\end{proof}
\begin{definition}[Mass Quantity]
  \label{def:physics:phyx-mini-0282:phyxminiproblems-problemphyxmini0282-massquantity}
  \lean{PhyXMiniProblems.ProblemPhyXMini0282.MassQuantity}
  A nonnegative physical mass
\end{definition}
\begin{proof}
  This is the declared model definition of Mass Quantity.
\end{proof}
\begin{definition}[Acceleration Quantity]
  \label{def:physics:phyx-mini-0282:phyxminiproblems-problemphyxmini0282-accelerationquantity}
  \lean{PhyXMiniProblems.ProblemPhyXMini0282.AccelerationQuantity}
  A nonnegative acceleration magnitude, of dimension L T⁻²
\end{definition}
\begin{proof}
  This is the declared model definition of Acceleration Quantity.
\end{proof}
\begin{definition}[Moment Of Inertia Quantity]
  \label{def:physics:phyx-mini-0282:phyxminiproblems-problemphyxmini0282-momentofinertiaquantity}
  \lean{PhyXMiniProblems.ProblemPhyXMini0282.MomentOfInertiaQuantity}
  A scalar moment of inertia, of dimension M L²
\end{definition}
\begin{proof}
  This is the declared model definition of Moment Of Inertia Quantity.
\end{proof}
\begin{definition}[Rotational Stiffness Quantity]
  \label{def:physics:phyx-mini-0282:phyxminiproblems-problemphyxmini0282-rotationalstiffnessquantity}
  \lean{PhyXMiniProblems.ProblemPhyXMini0282.RotationalStiffnessQuantity}
  The coefficient of angular displacement in a linearized restoring torque. Radians are dimensionless, so this has dimension M L² T⁻²
\end{definition}
\begin{proof}
  This is the declared model definition of Rotational Stiffness Quantity.
\end{proof}
\begin{definition}[length In Meters]
  \label{def:physics:phyx-mini-0282:phyxminiproblems-problemphyxmini0282-lengthinmeters}
  \lean{PhyXMiniProblems.ProblemPhyXMini0282.lengthInMeters}
  \uses{def:physics:phyx-mini-0282:phyxminiproblems-problemphyxmini0282-lengthquantity}
  Read a physical length in metres
\end{definition}
\begin{proof}
  This is the declared model definition of length In Meters.
\end{proof}
\begin{definition}[time In Seconds]
  \label{def:physics:phyx-mini-0282:phyxminiproblems-problemphyxmini0282-timeinseconds}
  \lean{PhyXMiniProblems.ProblemPhyXMini0282.timeInSeconds}
  \uses{def:physics:phyx-mini-0282:phyxminiproblems-problemphyxmini0282-timequantity}
  Read a physical duration in seconds
\end{definition}
\begin{proof}
  This is the declared model definition of time In Seconds.
\end{proof}
\begin{definition}[mass In Kilograms]
  \label{def:physics:phyx-mini-0282:phyxminiproblems-problemphyxmini0282-massinkilograms}
  \lean{PhyXMiniProblems.ProblemPhyXMini0282.massInKilograms}
  \uses{def:physics:phyx-mini-0282:phyxminiproblems-problemphyxmini0282-massquantity}
  Read a physical mass in kilograms
\end{definition}
\begin{proof}
  This is the declared model definition of mass In Kilograms.
\end{proof}
\begin{definition}[acceleration In Meters Per Second Squared]
  \label{def:physics:phyx-mini-0282:phyxminiproblems-problemphyxmini0282-accelerationinmeterspersecondsquared}
  \lean{PhyXMiniProblems.ProblemPhyXMini0282.accelerationInMetersPerSecondSquared}
  \uses{def:physics:phyx-mini-0282:phyxminiproblems-problemphyxmini0282-accelerationquantity}
  Read an acceleration in metres per second squared
\end{definition}
\begin{proof}
  This is the declared model definition of acceleration In Meters Per Second Squared.
\end{proof}
\begin{definition}[moment Of Inertia In Kilogram Meters Squared]
  \label{def:physics:phyx-mini-0282:phyxminiproblems-problemphyxmini0282-momentofinertiainkilogrammeterssquared}
  \lean{PhyXMiniProblems.ProblemPhyXMini0282.momentOfInertiaInKilogramMetersSquared}
  \uses{def:physics:phyx-mini-0282:phyxminiproblems-problemphyxmini0282-momentofinertiaquantity}
  Read a moment of inertia in kilogram metres squared
\end{definition}
\begin{proof}
  This is the declared model definition of moment Of Inertia In Kilogram Meters Squared.
\end{proof}
\begin{definition}[rotational Stiffness In Newton Meters]
  \label{def:physics:phyx-mini-0282:phyxminiproblems-problemphyxmini0282-rotationalstiffnessinnewtonmeters}
  \lean{PhyXMiniProblems.ProblemPhyXMini0282.rotationalStiffnessInNewtonMeters}
  \uses{def:physics:phyx-mini-0282:phyxminiproblems-problemphyxmini0282-rotationalstiffnessquantity}
  Read a linearized restoring coefficient in newton metres per radian
\end{definition}
\begin{proof}
  This is the declared model definition of rotational Stiffness In Newton Meters.
\end{proof}
\begin{definition}[Pivot]
  \label{def:physics:phyx-mini-0282:phyxminiproblems-problemphyxmini0282-pivot}
  \lean{PhyXMiniProblems.ProblemPhyXMini0282.Pivot}
  The two alternative suspension points named in the problem
\end{definition}
\begin{proof}
  This is the declared model definition of Pivot.
\end{proof}
\begin{definition}[Figure Point]
  \label{def:physics:phyx-mini-0282:phyxminiproblems-problemphyxmini0282-figurepoint}
  \lean{PhyXMiniProblems.ProblemPhyXMini0282.FigurePoint}
  Distinguished points drawn on the pendulum body
\end{definition}
\begin{proof}
  This is the declared model definition of Figure Point.
\end{proof}
\begin{definition}[Figure Label]
  \label{def:physics:phyx-mini-0282:phyxminiproblems-problemphyxmini0282-figurelabel}
  \lean{PhyXMiniProblems.ProblemPhyXMini0282.FigureLabel}
  Literal text labels visible in the primary image
\end{definition}
\begin{proof}
  This is the declared model definition of Figure Label.
\end{proof}
\begin{definition}[Pivot Mobility]
  \label{def:physics:phyx-mini-0282:phyxminiproblems-problemphyxmini0282-pivotmobility}
  \lean{PhyXMiniProblems.ProblemPhyXMini0282.PivotMobility}
  Mechanical role assigned to each candidate pivot
\end{definition}
\begin{proof}
  This is the declared model definition of Pivot Mobility.
\end{proof}
\begin{definition}[Axial Order]
  \label{def:physics:phyx-mini-0282:phyxminiproblems-problemphyxmini0282-axialorder}
  \lean{PhyXMiniProblems.ProblemPhyXMini0282.AxialOrder}
  Relative order of the three marked points along the pendulum body
\end{definition}
\begin{proof}
  This is the declared model definition of Axial Order.
\end{proof}
\begin{definition}[Body Treatment]
  \label{def:physics:phyx-mini-0282:phyxminiproblems-problemphyxmini0282-bodytreatment}
  \lean{PhyXMiniProblems.ProblemPhyXMini0282.BodyTreatment}
  The rigid-body approximation used for the pendulum
\end{definition}
\begin{proof}
  This is the declared model definition of Body Treatment.
\end{proof}
\begin{definition}[Oscillation Regime]
  \label{def:physics:phyx-mini-0282:phyxminiproblems-problemphyxmini0282-oscillationregime}
  \lean{PhyXMiniProblems.ProblemPhyXMini0282.OscillationRegime}
  The oscillation approximation under which the textbook period law holds
\end{definition}
\begin{proof}
  This is the declared model definition of Oscillation Regime.
\end{proof}
\begin{definition}[Reversible Pendulum Figure]
  \label{def:physics:phyx-mini-0282:phyxminiproblems-problemphyxmini0282-reversiblependulumfigure}
  \lean{PhyXMiniProblems.ProblemPhyXMini0282.ReversiblePendulumFigure}
  \uses{def:physics:phyx-mini-0282:phyxminiproblems-problemphyxmini0282-lengthquantity, def:physics:phyx-mini-0282:phyxminiproblems-problemphyxmini0282-pivot, def:physics:phyx-mini-0282:phyxminiproblems-problemphyxmini0282-figurepoint, def:physics:phyx-mini-0282:phyxminiproblems-problemphyxmini0282-figurelabel, def:physics:phyx-mini-0282:phyxminiproblems-problemphyxmini0282-pivotmobility, def:physics:phyx-mini-0282:phyxminiproblems-problemphyxmini0282-axialorder}
  Readout of the supplied image. Its distance fields are genuine physical lengths; Booleans record only visibility of graphical features
\end{definition}
\begin{proof}
  This is the declared model definition of Reversible Pendulum Figure.
\end{proof}
\begin{definition}[Reversible Physical Pendulum]
  \label{def:physics:phyx-mini-0282:phyxminiproblems-problemphyxmini0282-reversiblephysicalpendulum}
  \lean{PhyXMiniProblems.ProblemPhyXMini0282.ReversiblePhysicalPendulum}
  \uses{def:physics:phyx-mini-0282:phyxminiproblems-problemphyxmini0282-lengthquantity, def:physics:phyx-mini-0282:phyxminiproblems-problemphyxmini0282-timequantity, def:physics:phyx-mini-0282:phyxminiproblems-problemphyxmini0282-massquantity, def:physics:phyx-mini-0282:phyxminiproblems-problemphyxmini0282-accelerationquantity, def:physics:phyx-mini-0282:phyxminiproblems-problemphyxmini0282-momentofinertiaquantity, def:physics:phyx-mini-0282:phyxminiproblems-problemphyxmini0282-rotationalstiffnessquantity, def:physics:phyx-mini-0282:phyxminiproblems-problemphyxmini0282-pivot, def:physics:phyx-mini-0282:phyxminiproblems-problemphyxmini0282-bodytreatment, def:physics:phyx-mini-0282:phyxminiproblems-problemphyxmini0282-oscillationregime, def:physics:phyx-mini-0282:phyxminiproblems-problemphyxmini0282-reversiblependulumfigure}
  The single pendulum body, its two suspension configurations, and their observables. No field defines pivotSeparationL from the requested numerical answer
\end{definition}
\begin{proof}
  This is the declared model definition of Reversible Physical Pendulum.
\end{proof}
\begin{definition}[Matches Primary Figure]
  \label{def:physics:phyx-mini-0282:phyxminiproblems-problemphyxmini0282-matchesprimaryfigure}
  \lean{PhyXMiniProblems.ProblemPhyXMini0282.MatchesPrimaryFigure}
  \uses{def:physics:phyx-mini-0282:phyxminiproblems-problemphyxmini0282-lengthinmeters, def:physics:phyx-mini-0282:phyxminiproblems-problemphyxmini0282-reversiblephysicalpendulum}
  Geometric and qualitative facts visible in the primary image. In particular, the center of mass lies strictly closer to B than to A, so the equal-period pivots have unequal center-of-mass offsets. This excludes the degenerate midpoint case without assigning a numerical value to L
\end{definition}
\begin{proof}
  This is the declared model definition of Matches Primary Figure.
\end{proof}
\begin{definition}[Has Reversible Pendulum Problem Data]
  \label{def:physics:phyx-mini-0282:phyxminiproblems-problemphyxmini0282-hasreversiblependulumproblemdata}
  \lean{PhyXMiniProblems.ProblemPhyXMini0282.HasReversiblePendulumProblemData}
  \uses{def:physics:phyx-mini-0282:phyxminiproblems-problemphyxmini0282-lengthinmeters, def:physics:phyx-mini-0282:phyxminiproblems-problemphyxmini0282-timeinseconds, def:physics:phyx-mini-0282:phyxminiproblems-problemphyxmini0282-massinkilograms, def:physics:phyx-mini-0282:phyxminiproblems-problemphyxmini0282-accelerationinmeterspersecondsquared, def:physics:phyx-mini-0282:phyxminiproblems-problemphyxmini0282-momentofinertiainkilogrammeterssquared, def:physics:phyx-mini-0282:phyxminiproblems-problemphyxmini0282-reversiblephysicalpendulum}
  Textual data and standard textbook calibrations. Suspending from B restores the complete physical period observed at A; it does not assume any value for the required pivot separation
\end{definition}
\begin{proof}
  This is the declared model definition of Has Reversible Pendulum Problem Data.
\end{proof}
\begin{definition}[Satisfies Reversible Physical Pendulum Laws]
  \label{def:physics:phyx-mini-0282:phyxminiproblems-problemphyxmini0282-satisfiesreversiblephysicalpendulumlaws}
  \lean{PhyXMiniProblems.ProblemPhyXMini0282.SatisfiesReversiblePhysicalPendulumLaws}
  \uses{def:physics:phyx-mini-0282:phyxminiproblems-problemphyxmini0282-lengthinmeters, def:physics:phyx-mini-0282:phyxminiproblems-problemphyxmini0282-timeinseconds, def:physics:phyx-mini-0282:phyxminiproblems-problemphyxmini0282-massinkilograms, def:physics:phyx-mini-0282:phyxminiproblems-problemphyxmini0282-accelerationinmeterspersecondsquared, def:physics:phyx-mini-0282:phyxminiproblems-problemphyxmini0282-momentofinertiainkilogrammeterssquared, def:physics:phyx-mini-0282:phyxminiproblems-problemphyxmini0282-rotationalstiffnessinnewtonmeters, def:physics:phyx-mini-0282:phyxminiproblems-problemphyxmini0282-reversiblephysicalpendulum}
  The physical-pendulum laws at each suspension point: the parallel-axis theorem gives I\_p = I\_com + m d\_p²; gravity gives the linearized restoring coefficient m g d\_p; those two dimensionful quantities supply the generalized inertia and stiffness of Physlib's scalar harmonic oscillator; and the observed physical duration equals that oscillator's period. No field states the conjugate-pivot separation formula or a numerical answer
\end{definition}
\begin{proof}
  This is the declared model definition of Satisfies Reversible Physical Pendulum Laws.
\end{proof}
\begin{definition}[equivalent Simple Pendulum Length Meters]
  \label{def:physics:phyx-mini-0282:phyxminiproblems-problemphyxmini0282-equivalentsimplependulumlengthmeters}
  \lean{PhyXMiniProblems.ProblemPhyXMini0282.equivalentSimplePendulumLengthMeters}
  \uses{def:physics:phyx-mini-0282:phyxminiproblems-problemphyxmini0282-timeinseconds, def:physics:phyx-mini-0282:phyxminiproblems-problemphyxmini0282-accelerationinmeterspersecondsquared, def:physics:phyx-mini-0282:phyxminiproblems-problemphyxmini0282-reversiblephysicalpendulum}
  Equivalent simple-pendulum length computed from the period at A and local gravity. Its equality with the physical separation L remains a theorem
\end{definition}
\begin{proof}
  This is the declared model definition of equivalent Simple Pendulum Length Meters.
\end{proof}
\begin{definition}[Answer Choice]
  \label{def:physics:phyx-mini-0282:phyxminiproblems-problemphyxmini0282-answerchoice}
  \lean{PhyXMiniProblems.ProblemPhyXMini0282.AnswerChoice}
  Labels of the four displayed answer choices
\end{definition}
\begin{proof}
  This is the declared model definition of Answer Choice.
\end{proof}
\begin{definition}[answer Length Meters]
  \label{def:physics:phyx-mini-0282:phyxminiproblems-problemphyxmini0282-answerlengthmeters}
  \lean{PhyXMiniProblems.ProblemPhyXMini0282.answerLengthMeters}
  \uses{def:physics:phyx-mini-0282:phyxminiproblems-problemphyxmini0282-answerchoice}
  Distance in metres printed beside each answer choice
\end{definition}
\begin{proof}
  This is the declared model definition of answer Length Meters.
\end{proof}
\begin{definition}[recorded Answer Choice]
  \label{def:physics:phyx-mini-0282:phyxminiproblems-problemphyxmini0282-recordedanswerchoice}
  \lean{PhyXMiniProblems.ProblemPhyXMini0282.recordedAnswerChoice}
  \uses{def:physics:phyx-mini-0282:phyxminiproblems-problemphyxmini0282-answerchoice}
  The answer label recorded in the supplied dataset metadata
\end{definition}
\begin{proof}
  This is the declared model definition of recorded Answer Choice.
\end{proof}
\begin{definition}[Is Unique Closest Answer Choice]
  \label{def:physics:phyx-mini-0282:phyxminiproblems-problemphyxmini0282-isuniqueclosestanswerchoice}
  \lean{PhyXMiniProblems.ProblemPhyXMini0282.IsUniqueClosestAnswerChoice}
  \uses{def:physics:phyx-mini-0282:phyxminiproblems-problemphyxmini0282-answerchoice, def:physics:phyx-mini-0282:phyxminiproblems-problemphyxmini0282-answerlengthmeters}
  A displayed choice is strictly closer than every distinct alternative
\end{definition}
\begin{proof}
  This is the declared model definition of Is Unique Closest Answer Choice.
\end{proof}
% --- Archon declaration topology end ---
% --- Archon physics formalization source end ---
